\chapter{Összegzés}

A féléves projektben elkészítettem egy AES-alapú titkosító rendszert, amely:

\begin{itemize}
  \item több üzemmódot kezel (CTR, GCM; natív oldalon CBC is elérhető),
  \item több végrehajtási utat kínál (natív CPU, OpenCL, managed .NET referencia),
  \item és egy WinForms felületen keresztül mérhető/összehasonlítható.
\end{itemize}

\section{Válaszok az általam vizsgált kérdésekre}

\subsection{x64 vagy x86 gyorsabb?}

A mérések alapján:

\begin{itemize}
  \item \textbf{Natív CPU}: jellemzően \textbf{x64 gyorsabb} (PC2/PC3/PC4).
  \item \textbf{Managed .NET}: szintén \textbf{x64 gyorsabb}, sok esetben nagy különbséggel.
  \item \textbf{OpenCL}: a throughput \textbf{nagyon hasonló} x86 és x64 között; ahol diszkrét GPU dolgozik (PC4),
        gyakorlatilag 1\% körüli eltéréseket látni.
\end{itemize}

Ennek alapján én azt mondom, hogy az x64 előnye főleg a CPU/JIT oldalon jön ki,
az OpenCL kernel futásideje viszont kevésbé érzékeny a host architektúrára.

\subsection{Befolyásolja-e a töltő (PC3) a gyorsaságot?}

Igen. PC3-on töltőn minden engine gyorsabb, különösen a CPU-intenzív esetek.
Ezt a power management (órajel/boost limit, energiatakarékossági profil) természetes következményének tartom.

\section{További megfigyelések}

\begin{itemize}
  \item \textbf{Kulcsméret növelése} általában csökkenti a throughputot (CTR és GCM esetén is).
  \item \textbf{Nagyobb bemenet} OpenCL-nél segít: az overhead jobban amortizálódik, ezért 256 MB-nál jellemzően jobb throughputot mértem, mint 64 MB-nál.
  \item \textbf{PC1 OpenCL x86} futtatás nem sikerült: a legvalószínűbb ok a 32 bites OpenCL runtime/driver hiánya.
\end{itemize}

\section{Lehetséges fejlesztési irányok}

Ha tovább vinném a projektet, a következőket tartanám a legértékesebbnek:

\begin{itemize}
  \item Natív CPU oldalon AES-NI / vektorizálás (SIMD) bevezetése, hogy a baseline ne legyen ennyire távol a modern implementációktól.
  \item OpenCL oldalon a host--device adatmozgatás csökkentése (pl. pinned memory, nagyobb chunk-ok, overlapped copy/compute).
  \item GCM GHASH rész további párhuzamosítása (pl. nagyobb blokkosítás, precomputált táblák, jobb redukciós stratégia).
\end{itemize}

\cleardoublepage
